\subsection{A mixed-parity divisibility obstruction}
\label{sec:mixed-parity-congruence}

The distribution of hypothetical integer roots modulo four gives further
infinite classes with unequal exponent valuations.

\begin{lemma}
\label{lem:three-odd-roots}
Let $f\in\mathbb Z[x]$ be monic of degree five with
$f(x)\equiv x^2(x+1)^3\pmod2$. If all its roots are integers,
then $32\mid f(1)$ or $32\mid f(3)$.
\end{lemma}

\begin{proof}
The binary factorization forces precisely three odd roots and two even
roots, counted with multiplicity. At least two odd roots share a residue
$r\in\{1,3\}$ modulo four. In $f(r)$, their two factors are divisible
by four and the third odd-root factor is divisible by two. The even-root
factors are odd. Therefore $32\mid f(r)$, including when $f(r)=0$.
\end{proof}

\begin{theorem}
\label{thm:mixed-parity-congruence}
Let $a,b$ be the two odd exponents and $c$ the even exponent of a positive
three-prime exponent triple. Each row in the following table, including
exchange of $a,b$, implies that its ideal intersection graph is nonintegral.
\begin{center}
\begin{tabular}{ccc|cc}
\toprule
$a\bmod8$ & $b\bmod8$ & $c\bmod8$ & $h_C(1)\bmod32$ & $h_C(3)\bmod32$\\
\midrule
1&3&2&8&16\\
1&7&2&24&16\\
3&7&2&16&8\\
3&5&6&8&16\\
3&7&6&16&24\\
5&7&6&24&16\\
\bottomrule
\end{tabular}
\end{center}
\end{theorem}

\begin{proof}
For either mixed exponent parity,
$h_C(x)\equiv x^2(x+1)^3\pmod2$. Substituting
$(a,b,c)=(8u+r_1,8v+r_2,8w+r_3)$ in the general quotient quintic gives
the last two columns: every nonconstant coefficient in $u,v,w$ is
divisible by $32$. Both entries are nonzero modulo $32$, contradicting
Lemma~\ref{lem:three-odd-roots} if all quotient roots were integers.
The resulting noninteger real complement eigenvalue has a noninteger
graph lift.
\end{proof}

For example, $(10,17,19)$ and $(14,19,21)$ meet these criteria with
gcd one and unequal $2$-adic valuations. There is no minimum or ratio
bound in the argument. The weaker divisibilities $4\mid h_C(2t)$
and $8\mid h_C(2t+1)$ hold identically for both mixed exponent parities;
they alone do not supply the contradiction. The distribution of the
three odd roots modulo four is the additional constraint.
